\documentclass{beamer}
\usepackage[utf8]{inputenc}
\usepackage[T1]{fontenc}
\usepackage{amsmath}
\usepackage{amssymb}
\usepackage{graphicx}
\usepackage{longtable}
\usepackage{array}
\usepackage{adjustbox}
\usepackage{ragged2e}
\usepackage[colorlinks=true, linkcolor=black, urlcolor=black, citecolor=black]{hyperref}

% Theme setup
\usetheme{CambridgeUS}
\usecolortheme{dolphin}
\setbeamertemplate{navigation symbols}{}

% Header & Footer (no text)
\setbeamertemplate{headline}{
  \begin{beamercolorbox}[wd=\paperwidth,ht=2.5ex,dp=1.125ex]{section in head/foot}%
  \end{beamercolorbox}%
  \begin{beamercolorbox}[wd=\paperwidth,ht=2.5ex,dp=1.125ex]{subsection in head/foot}%
  \end{beamercolorbox}%
}

\setbeamertemplate{footline}{
  \leavevmode%
  \hbox{%
  \begin{beamercolorbox}[wd=\paperwidth,ht=2.5ex,dp=1.125ex,center]{author in head/foot}%
  \end{beamercolorbox}}%
  \vskip0pt%
}

% Fonts
\setbeamerfont{title}{size=\Large}
\setbeamerfont{subtitle}{size=\large}
\setbeamerfont{author}{size=\normalsize}
\setbeamerfont{institute}{size=\normalsize}
\setbeamerfont{date}{size=\small}
\setbeamerfont{frametitle}{size=\Large}

\title{Apricity : Startup Success Predictor using Multi-Agent System}
\author[Group 8]{
\textbf{Group 8}\\[6pt]
\begin{tabular}{ll}
Darwin Thomas & 22CSB13 MDL22CS065 \\
K V Varun Krishnan & 22CSB30 MDL22CS114 \\
Manu S Panicker & 22CSB31 MDL22CS119 \\
M Arjun Krishna & 22CSB32 MDL22CS122
\end{tabular}
}
\institute{\textbf{Model Engineering College}}
\date{October 2025}

\begin{document}

% ---------------- Title Page ----------------
\begin{frame}
  \titlepage
\end{frame}

% ---------------- Table of Contents ----------------
\begin{frame}{Contents}
\tableofcontents[sections={1-8}]
\end{frame}

\begin{frame}{Contents}
\tableofcontents[sections={9-16}]
\end{frame}

% ======================== INTRODUCTION ========================
\section{Introduction}

\begin{frame}
\centering
\Huge \textbf{Introduction}
\end{frame}

\begin{frame}{Introduction}
\begin{itemize}
\item Apricity aims to assist founders and investors by providing fast, data-driven evaluation of startup viability
\item Designed as an AI-powered software prototype that simulates expert-level startup assessment through autonomous agents
\item Employs a multi-agent architecture where each agent represents a domain expert such as Finance, Technology, Marketing, Sales, and Product
\item Integrates Large Language Models with Retrieval-Augmented Generation (RAG) to ensure informed, context-aware, and explainable predictions
\item Generates structured insights including risk analysis, strengths, weaknesses, and improvement suggestions in real time

\end{itemize}
\end{frame}


% ======================== FEATURES ========================
\section{Features}

\begin{frame}
\centering
\Huge \textbf{Features}
\end{frame}

\begin{frame}{Key Features of the System}
\begin{itemize}
    \item \textbf{AI-Powered Evaluation:} Uses large language models to assess startup viability across multiple dimensions.
    \item \textbf{Multi-Agent Architecture:} Independent agents simulate domain experts such as Financial Analyst, VC, CTO, Marketing, Sales, and Product specialists.
    \item \textbf{Retrieval-Augmented Generation (RAG):} Grounds evaluations using up-to-date external knowledge sources for improved accuracy.
    
    \item \textbf{Real-Time Feedback:} Generates comprehensive startup assessments within minutes.
    \item \textbf{Modular and Scalable Design:} Agents can be added, removed, or customized for different industries and evaluation criteria.
\end{itemize}
\end{frame}


% ======================== PURPOSE ========================
\section{Purpose and Scope}

\begin{frame}
\centering
\Huge \textbf{Purpose and Scope}
\end{frame}

\begin{frame}{Purpose and Scope}
\begin{itemize}
    \item Defines the design and implementation of an AI-driven system for predicting startup success
    \item Captures both functional and non-functional requirements of the Startup Success Predictor software prototype
    \item Simulates expert-driven startup evaluation using autonomous multi-agent reasoning
    \item Focuses on software components including agent orchestration, retrieval pipeline, and evaluation logic
    \item Provides structured, explainable assessments across finance, technology, market, sales, and product domains
    \item Designed to be extensible for future integration with real-world datasets, investor tools, and decision-support platforms
\end{itemize}
\end{frame}



% ======================== PROBLEM STATEMENT ========================
\section{Problem Statement}

\begin{frame}
\centering
\Huge \textbf{Problem Statement}
\end{frame}

\begin{frame}{Problem Statement}
\begin{itemize}
    \item Startup evaluation today relies heavily on manual expert reviews.
    \item These processes are slow, expensive, and difficult to scale.
    \item Early-stage founders often lack access to structured, expert-level feedback.
\end{itemize}
\end{frame}

% ======================== PROPOSED SOLUTION ========================
\section{Proposed Solution}

\begin{frame}
\centering
\Huge \textbf{Proposed Solution}
\end{frame}

\begin{frame}{Proposed Solution}
\begin{itemize}
    \item An AI-driven multi-agent system that simulates expert startup evaluation.
    \item Each agent represents a domain expert (finance, product, market, etc.).
    \item Produces structured insights instead of a single opaque prediction.
\end{itemize}
\end{frame}

% ======================== PROGRESS OVERVIEW ========================
% ======================== UPDATED PROGRESS OVERVIEW ========================
\section{Project Progress Overview}

\begin{frame}
\centering
\Huge \textbf{Project Progress Overview}
\end{frame}

\begin{frame}{Project Progress Overview (Current)}
\begin{itemize}
    \item Overall Status: \textbf{$\sim$90\%} core MVP pipeline implemented
    \item Frontend: Guided startup workspace (templates + rich-text editor), autosave, PDF upload + viewer
    \item Backend: FastAPI endpoint runs \textbf{multi-agent analysis} with RAG chunking + retrieval
    \item Web enrichment: \textbf{Exa API} competitor discovery implemented (structured competitor list)
    \item Report: Results rendered in UI + \textbf{PDF export implemented} via \textbf{html2pdf.js}
\end{itemize}
\end{frame}

% ======================== SYSTEM FLOW DIAGRAM ========================
\section{System Flow}

\begin{frame}
\centering
\Huge \textbf{System Flow Diagram}
\end{frame}

\begin{frame}{End-to-End System Flow}
\centering
\includegraphics[
    width=0.9\textwidth,
    height=0.75\textheight,
    keepaspectratio
]{Block.png}

\vspace{0.5em}
\small Flow diagram showing end-to-end data movement across system modules.
\end{frame}


% ======================== FLOW EXPLANATION ========================

\begin{frame}{System Flow Explanation (1/2)}
\small
\begin{enumerate}
    \item \textbf{User Input → Workspace (Frontend):}  
    User fills guided startup sections using artifact templates (rich-text editor) and can upload a PDF for reference.

    \item \textbf{Workspace → Local Persistence:}  
    Documents, files, and view state are autosaved to \textbf{IndexedDB} (debounced) for reliability across refreshes.

    \item \textbf{Workspace → “Develop a View”:}  
    Frontend aggregates all document content into a single analysis prompt and creates a new \texttt{viewId} with \texttt{pending} status.

    \item \textbf{View Page → Analysis Request:}  
    UI shows progress stages and sends the aggregated prompt to the internal API endpoint \texttt{/api/views}.

    \item \textbf{API Route → Backend Forwarding Setup:}  
    Next.js API validates the request and forwards it to the Python backend endpoint \texttt{/view}.
\end{enumerate}
\end{frame}

\begin{frame}{System Flow Explanation (2/2)}
\small
\setcounter{enumi}{5}
\begin{enumerate}
    \item \textbf{Next.js API → FastAPI Backend:}  
    The prompt is posted to \texttt{POST /view} for analysis.

    \item \textbf{Backend → RAG Context Preparation:}  
    Input is chunked, embedded (OpenAI embeddings), and indexed in \textbf{Chroma} to enable retrieval during agent reasoning.

    \item \textbf{Backend → Multi-Agent Evaluation:}  
    \textbf{LangGraph} runs domain agents (Finance, VC, CTO, Marketing, Product) in parallel the shared retriever.

    \item \textbf{Backend → Web Enrichment (Competitors):}  
    \textbf{Exa API} is used to discover and enrich top competitors into a structured competitor dataset (disabled if \texttt{EXA\_API\_KEY} is missing).

    \item \textbf{Backend → Structured JSON Response:}  
    Backend returns agent analyses + competitor search status + competitor list as JSON.

    \item \textbf{Frontend → Results Rendering \& Storage:}  
    The view is marked \texttt{completed} (or \texttt{error}), results are rendered section-wise, and stored in the workspace state for later access.

    \item \textbf{Frontend → Report Export:}  
    User exports a styled A4 PDF report using \textbf{html2pdf.js} (markdown rendering supported for agent outputs).
\end{enumerate}
\end{frame}


% ======================== INPUT MODULE ========================
\section{Input Module (React / Next.js)}

\begin{frame}
\centering
\Huge \textbf{Input Module}
\end{frame}

\begin{frame}{Input Module – Planned Functionality}
\begin{itemize}
    \item Collect startup inputs (problem, solution, market, team, funding, risks)
    \item Support document uploads (pitch/report) as additional context
    \item Normalize and prepare input into a single structured context for analysis
\end{itemize}
\end{frame}

\begin{frame}{Input Module – Current Implementation}
\begin{itemize}
    \item Implemented a \textbf{workspace-style UI} (Next.js + React)
    \item Rich-text editor implemented using \textbf{TipTap} with guided templates (artifacts)
    \item Dynamic placeholder prompts help users fill required sections
    \item \textbf{Autosave implemented} using IndexedDB (\texttt{idb-keyval})
    \item PDF upload supported + PDF viewer (stored locally for reference)
    \item “Develop a View” aggregates all document text and sends it for backend analysis
\end{itemize}

\vspace{0.4em}
\textbf{Status:} \textbf{90\%}
\end{frame}

\begin{frame}{Input Module – Pending Work}
\begin{itemize}
    \item Improve UI/UX polish (navigation, empty states, guidance clarity)
\end{itemize}
\end{frame}

\begin{frame}{Module Demonstration }
\centering
\IfFileExists{first.png}{
\includegraphics[
    width=0.85\textwidth,
    height=0.65\textheight,
    keepaspectratio
]{first.png}

\vspace{0.5em}
\textbf{Description:} Workspace editor + guided templates.
}{
\fbox{\parbox{0.85\textwidth}{Screenshot placeholder: \texttt{Screenshots/SSP/workspace.png}}}
}
\end{frame}


% ======================== WEB SEARCHER MODULE ========================
\section{Web Searcher Module (Exa API)}

\begin{frame}
\centering
\Huge \textbf{Web Searcher Module}
\end{frame}

\begin{frame}{Web Searcher Module – Planned Functionality}
\begin{itemize}
    \item Enrich startup context using up-to-date web sources
    \item Fetch competitors, company facts, and market signals
    \item Merge web findings into a structured enrichment dataset
\end{itemize}
\end{frame}

\begin{frame}{Web Searcher Module – Current Implementation}
\begin{itemize}
    \item Competitor enrichment implemented using \textbf{Exa Websets API}
    \item The system generates a compact search sentence (LLM-based) for better retrieval
    \item Outputs a \textbf{structured competitor list} including:
    \begin{itemize}
        \item Company name, website, description
        \item Funding, revenue (if available), employee count
        \item Founded year, target customer, headquarters
    \end{itemize}
    \item Social-signal collector implemented as a \textbf{CLI service} (Reddit + Hacker News + optional X), currently \textbf{not wired} into the main API
\end{itemize}

\vspace{0.4em}
\textbf{Status:} \textbf{90\%}
\end{frame}

\begin{frame}{Web Searcher Module – Pending Work}
\begin{itemize}
    \item Add \textbf{source attribution} (URLs / evidence pointers) into enrichment output
    \item Add caching + rate-limit handling for API stability
    \item Integrate social-signal output into backend response (optional feature toggle)
    \item Improve enrichment quality filtering (dedupe + conflict handling)
\end{itemize}
\end{frame}


% ======================== AGENT MODULE ========================
\section{Agent Module (LangGraph + OpenAI)}

\begin{frame}
\centering
\Huge \textbf{Agent Module}
\end{frame}

\begin{frame}{Agent Module – Planned Functionality}
\begin{itemize}
    \item Run expert-style evaluations across multiple domains
    \item Produce concise reasoning per domain (finance, product, tech, go-to-market, VC)
    \item Keep outputs structured for report generation
\end{itemize}
\end{frame}

\begin{frame}{Agent Module – Current Implementation}
\begin{itemize}
    \item Multi-agent analysis implemented using \textbf{LangGraph} + \textbf{OpenAI Chat models}
    \item Implemented agents:
    \begin{itemize}
        \item Financial Analyst
        \item Venture Capitalist
        \item CTO
        \item Marketing Specialist
        \item Product Analyst
    \end{itemize}
    \item Implemented RAG pipeline:
    \begin{itemize}
        \item Chunking via RecursiveCharacterTextSplitter
        \item Retrieval via \textbf{Chroma} + OpenAI embeddings
    \end{itemize}
    \item Parallel execution implemented (async orchestration + aggregation)
\end{itemize}

\vspace{0.4em}
\textbf{Status:} \textbf{90\%}
\end{frame}

\begin{frame}{Agent Module – Pending Work}
\begin{itemize}
    \item Prompt refinement for higher consistency and clearer outputs
    \item Token optimization (shorter context + better chunk selection)
    \item Optional: add missing agents (e.g., Sales) and unify output schema
\end{itemize}
\end{frame}


\begin{frame}{Module Demonstration }
\centering
\IfFileExists{loading.png}{
\includegraphics[
    width=0.85\textwidth,
    height=0.65\textheight,
    keepaspectratio
]{loading.png}

\vspace{0.5em}
\textbf{Description:} Multi Agent answer loading.
}{
\fbox{\parbox{0.85\textwidth}{Screenshot placeholder: \texttt{Screenshots/SSP/workspace.png}}}
}
\end{frame}


\begin{frame}{Module Demonstration }
\centering
\IfFileExists{result.png}{
\includegraphics[
    width=0.85\textwidth,
    height=0.65\textheight,
    keepaspectratio
]{result.png}

\vspace{0.5em}
\textbf{Description:} Multi Agent Analysis.
}{
\fbox{\parbox{0.85\textwidth}{Screenshot placeholder: \texttt{Screenshots/SSP/workspace.png}}}
}
\end{frame}

% ======================== REPORT GENERATION MODULE ========================
\section{Report Generation Module (UI + PDF Export)}

\begin{frame}
\centering
\Huge \textbf{Report Generation Module}
\end{frame}

\begin{frame}{Report Generation Module – Planned Functionality}
\begin{itemize}
    \item Present final evaluation in a readable, structured report
    \item Provide exportable output (PDF) for sharing
\end{itemize}
\end{frame}

\begin{frame}{Report Generation Module – Current Implementation}
\begin{itemize}
    \item Implemented results page that renders backend outputs per section
    \item Implemented PDF export in the frontend using \textbf{html2pdf.js}
    \item PDF styling includes header/metadata + section formatting (markdown rendering supported)
    \item Progress UI implemented for “view generation” (stage indicator + error handling)
\end{itemize}

\vspace{0.4em}
\textbf{Status:} \textbf{90\%}
\end{frame}

\begin{frame}{Report Generation Module – Pending Work}
\begin{itemize}
    \item More polished PDF layout (tables for competitors, consistent spacing, page breaks)
    \item Add higher-level summary blocks (highlights, risks, recommendations)
    \item Optional: add computed scoring + rubric (not currently implemented)
\end{itemize}
\end{frame}

\begin{frame}{Module Demonstration (Optional Screenshot)}
\centering
\IfFileExists{pdf.png}{
\includegraphics[
    width=0.85\textwidth,
    height=0.65\textheight,
    keepaspectratio
]{pdf.png}

\vspace{0.5em}
\textbf{Description:} PDF export of analysis report.
}{
\fbox{\parbox{0.85\textwidth}{Screenshot placeholder: \texttt{Screenshots/SSP/pdf-export.png}}}
}
\end{frame}



% ======================== CONCLUSION ========================
% ======================== PROJECT SUMMARY ========================
\section{Progress Summary}

\begin{frame}{Conclusion \& Next Steps}
\begin{itemize}
    \item Project is on track with a \textbf{modular, incremental implementation} toward an MVP
    \item Core modules are \textbf{implemented and working} end-to-end:
    \begin{itemize}
        \item Input workspace (\textbf{Next.js + React}): guided templates, rich-text editor, PDF upload/viewer, autosave
        \item Web enrichment (\textbf{Exa API}): competitor discovery + structured competitor dataset
        \item Multi-agent evaluation (\textbf{LangGraph + OpenAI}): domain agents with RAG-based context retrieval
        \item Report output: results UI + \textbf{PDF export} implemented (html2pdf.js)
    \end{itemize}
    \item Current focus: \textbf{performance and quality optimization} (prompt tuning, token efficiency, enrichment reliability, UI/UX polish)
    \item Next milestone: improved citations/source attribution + caching/rate limits + more polished, consistent PDF reports
\end{itemize}
\end{frame}



% ======================== REFERENCES ========================

\section{References}

\begin{frame}
\centering
\Huge \textbf{References}
\end{frame}

\begin{frame}
\frametitle{References}
\footnotesize{

[1] A. Vaswani, N. Shazeer, N. Parmar, et al., “Attention Is All You Need,” in \textit{Advances in Neural Information Processing Systems}, 2017. \\

[2] M. Puvvadi, S. K. Arava, A. Santoria, et al., “Survey of Marketing Agents, Agent-Based Models, and Generative AI in Marketing,” \textit{arXiv preprint} arXiv:2401.01234, 2024. \\

[3] S. Han, Q. Zhang, Y. Yao, W. Jin, and Z. Xu, “LLM Multi-Agent Systems: Challenges and Open Problems,” \textit{arXiv preprint} arXiv:2403.04567, 2024. \\

[4] E. Gavrilenko, F. Khosmood, M. Rastad, et al., “Improving Startup Success with Text Analysis,” \textit{JP Morgan Chase Technical Report}, 2023. \\

[5] A. Khan and M. V. Mäntylä, “Developing RAG-Based LLM Systems from PDFs,” \textit{Tampere University}, 2024. \\
}
\end{frame}

\begin{frame}
\frametitle{References}
\footnotesize{
[6] A. Yehudai, L. Eden, A. Li, et al., “Survey on Evaluation of LLM-Based Agents,” \textit{arXiv preprint} arXiv:2402.07890, 2024. \\

[7] S. Xiong and Y. Ihlamur, “Founder-GPT: Self-Play to Evaluate the Founder–Idea Fit,” \textit{arXiv preprint} arXiv:2404.01111, 2024. \\

[8] G. Zhang, L. Niu, et al., “MaAS – Multi-Agent Architecture Search via Agentic Supernet,” in \textit{Proceedings of ICML}, 2025. \\

[9] P. Yin, J. Liu, et al., “Zero-Indexing: Decoupling Planning and Execution in LLMs,” Stanford, UC Berkeley, Microsoft Research, 2024. \\

[10] K. Lee, J. Chen, et al., “Chain of Agents: Modular Composition for Complex Reasoning,” \textit{Stanford University}, 2024. \\

[11] X. Wang, Y. Ihlamur, and F. Alican, “SSFF: Investigating LLM Predictive Capabilities for Startup Success,” \textit{arXiv preprint} arXiv:2405.02777, 2024. \\

}
\end{frame}

% ======================== THANK YOU ========================
\begin{frame}
\centering
\Huge \textbf{Thank You}
\end{frame}

\end{document}